\documentclass[10pt,a4paper]{amsart}
\usepackage[margin=19mm]{geometry}
\usepackage{amsmath,amssymb,mathtools}
\usepackage[scr=rsfs,cal=euler]{mathalfa}
\usepackage{xfrac}
\usepackage[unicode,hidelinks,bookmarks=false]{hyperref}
\newcommand{\ie}{i.e.\ }
\newcommand{\cf}{cf.\ }
\newcommand{\resp}{respectively\ }
\newcommand{\Subsection}{\subsection}
\providecommand{\nobreakdash}{\nobreak\hbox{-}}
\setlength{\emergencystretch}{2em}
\multlinegap0pt
\usepackage{booktabs}
\usepackage{nicefrac}
\usepackage{xcolor}
\usepackage[shortlabels]{enumitem}
\usepackage{multirow}
\usepackage{amssymb}
\usepackage{mathtools}
\newtheorem{theorem}{Theorem}
\usepackage{cleveref}
\crefname{theorem}{Theorem}{Theorems}
\newcommand{\I}{\mathcal{I}}
 \newcommand{\norm}[1]{{\left\|{#1}\right\|}}
\title{Accelerated Projected Gradient Algorithms for Sparsity Constrained Optimization Problems}
\author{Jan Harold Alcantara, Ching-pei Lee}
\date{Source: arXiv:2211.02271v2 (CC BY 4.0)}
\begin{document}
\maketitle

\noindent\emph{Source and licence.} Accelerated Projected Gradient Algorithms for Sparsity Constrained Optimization Problems. Version arXiv:2211.02271v2 (CC BY 4.0), distributed by arXiv under the Creative Commons Attribution 4.0 International licence, http://creativecommons.org/licenses/by/4.0/, as recorded in the arXiv licence metadata for this exact version and shown on the abstract page https://arxiv.org/abs/2211.02271v2. Full licence notice: \texttt{licenses/arxiv-2211.02271-CC-BY-4.0.txt}.\par\smallskip

\noindent\textit{Editorial scope.} Counted unit: the article's theorem thm:superlinear (source0.tex lines 4361--4382). The article's complete original proof of it is delivered with it (source0.tex lines 4383--4406, one \texttt{proof} environment of 24 lines). No mathematics is written, repaired or re-derived: the statement and the proof are the article's own lines, in the article's own notation, and no retained line is substituted or re-flowed.  Retained uncounted as the source-local context the proof needs: lines 212--215.  Nothing was omitted from any retained span, so the package carries no omission record.  No retained line contains a citation, so the wrapper carries no bibliography and no bibliography entry.  No retained line contains a figure environment, an image-inclusion command or drawing code, so no image is delivered and the package declares no asset dependency.  Editorial formatting only, in addition: an AMS \texttt{amsart} preamble that restates the article's own declarations and macros used by the retained lines, namely the environments \texttt{theorem} on separate counters, together with the standard packages those lines need.  The retained environments print in this document as Theorem 1 in order of appearance, whereas the article prints the same items with the numbering of its own full document.  The complete native source of the article as distributed in the arXiv e-print for this exact version is bundled unchanged as \texttt{sources/133005/source0.tex}.
	\begin{equation}
	\label{eq:problem}
	\min\nolimits _{w \in A_s} f(w),
	\end{equation}

\begin{theorem}
\label{thm:superlinear}
Assume that we have a mapping $T_1(w)$ such that its generated
iterates $\{w^k\}$ with $w^{k+1} \in T_1(w^k)$ converge to a
stationary point $w^*$ of \cref{eq:problem} and
\begin{equation}
	\norm{\hat w - T_1(w^*)} \leq \norm{w - w^*}, \quad
	\forall \hat w \in T_1 \left( w \right)
	\label{eq:nonexp}
\end{equation}
for all $w$ in a neighborhood $U$ of $w^*$ and in some $A_J$ with $J$
satisfying $J \in \I_{w^*}$,
and that there is another mapping $T_2$ that, when given an initial
point $w^0 \in A_J$, generates iterates that are all in $A_J$ and
superlinearly convergent to $w^*$ within $U$ for each $J \in
\I_{w^*}$ with $T_2(w^*) = w^*$, then the iterates generated by
\begin{equation}
	w^{k+1} \in T_2\left(T_1\left(w^{k}\right)\right)
	\label{eq:seq}
\end{equation}
converge to $w^*$ at the same superlinear rate as that of $T_2$.
\end{theorem}

\begin{proof}
We assume without loss of generality that
\begin{equation}
	\norm{T_2(w) - w^*} \leq c \norm{w - w^*}^{1+\rho}
	\label{eq:superlinear}
\end{equation}
for some $c, \rho > 0$ for all $w \in A_J \cap U$ for all $J \in
\I_{w^*}$.
Then by \cref{eq:nonexp}, and by denoting
\[
	\hat w^{k+1} \in T_1(w^k),
\]
as the element in $T_1(w^k)$ leading to $w^{k+1}$,
we obtain
\begin{align*}
	\norm{w^{k+1} - w^*} &= \norm{T_2 \left(\hat w^k  \right) -
	w^*}\\
	&\leq  c \norm{\hat w^k - w^*}^{1+\rho}\\
	&= c \norm{T_1(w^k) - w^*}^{1+\rho} \\
	&= c \norm{w^k - w^*}^{1+\rho},
\end{align*}
where the the first inequality is from \cref{eq:superlinear}.
Therefore, the conclusion of the theorem is proven.	
\end{proof}

\end{document}
